\documentclass[10pt]{article}

\usepackage[margin=1in]{geometry}
\usepackage{amsmath,amssymb,bm,mathtools}
\usepackage{microtype}
\usepackage[hyphens]{url}
\usepackage{hyperref}
\urlstyle{same}

\usepackage{algorithm}
\usepackage{algpseudocode}
\usepackage{float} % for [H]

% ---------- Notation helpers ----------
\DeclareMathOperator*{\argmin}{arg\,min}
\newcommand{\Map}{\mathrm{Map}}
\newcommand{\MOCU}{\mathrm{MOCU}}
\newcommand{\Safe}{\mathrm{Safe}}
\newcommand{\wmed}{\mathrm{wmed}}
\newcommand{\Ninner}{N_{\mathrm{inner}}}
\newcommand{\Nouter}{N_{\mathrm{outer}}}
\newcommand{\ninner}{n_{\mathrm{inner}}}
\newcommand{\nouter}{n_{\mathrm{outer}}}

% ---------- Algorithmic style ----------
\algrenewcommand\algorithmicrequire{\textbf{Input:}}
\algrenewcommand\algorithmicensure{\textbf{Output:}}
\algrenewcommand\algorithmicreturn{\textbf{Return}}

\begin{document}
\pagestyle{plain}

\title{\textbf{DAD--MOCU for Power-System Swing Dynamics:\\
Sequential Bayesian Experimental Design for Minimizing Decision-Relevant Uncertainty}}
\author{Gaoming Lin \;\; (Advisor: Dr.\ Byung-Jun Yoon)\\
Department of Electrical and Computer Engineering, Texas A\&M University}
\date{January 2026}

\maketitle

\vspace{-0.5em}
\noindent
\textbf{Summary.}
This document gives the full pseudocode and mathematical setup for the \textbf{DAD--MOCU} framework: \textbf{sequential Bayesian optimal experimental design (sBOED)} on power-system \textbf{swing (second-order Kuramoto)} dynamics over an IEEE-14 bus network. The pipeline is: choose a \textbf{probe} (design $\xi$) $\to$ run the system and get an \textbf{observation} $y$ (max ROCOF + noise) $\to$ update the \textbf{belief} over parameters via \textbf{Bayes' rule} $\to$ compute the \textbf{Bayes-optimal control gain} $\hat{\gamma}(p)$ and \textbf{MOCU} (mean objective cost of uncertainty). The \textbf{decision objective} is the \emph{minimal safe gain} $\gamma^\ast(\vartheta)$: the smallest gain such that the system stays within frequency and ROCOF limits under a fixed reference contingency; $\gamma^\ast$ is evaluated under that reference probe and candidate control. The \emph{target} of the sequential design is to minimize \emph{expected terminal MOCU} at step $t=T$, i.e.\ $\mathbb{E}[\mathrm{MOCU}(p_T)]$. Belief index $t$ runs from $0$ (prior $p_0$; all methods share the same MOCU here) to $T$ (belief $p_T$ after $T$ steps); we apply designs at steps $t=1,\dots,T$. A \emph{myopic} (greedy) rule that minimizes only one-step-ahead expected MOCU is \emph{generally suboptimal} for terminal MOCU; we use an \emph{amortized} policy (neural network $\pi_\phi$, trained offline) to select designs.

\medskip
\noindent
\textbf{Reader's guide.}
\begin{itemize}
\item \textbf{\S\ref{sec:setup} (Framework):} Defines the goal, dynamics, latent parameters, prior, design space, observation model (ROCOF), likelihood, posterior update (including log-domain for numerical stability), cost $J$, Bayes-optimal gain $\hat{\gamma}(p)$, MOCU, and the optimization objective. All notation and assumptions needed for the algorithms are here.
\item \textbf{\S\ref{sec:sequential} (Algorithm 1):} One \emph{episode} of sequential design using a \emph{pre-trained} amortized policy $\pi_\phi$: at each step we sample a probe from the policy, observe $y_t$, update the belief in log-domain, then compute $\hat{\gamma}(p_t)$ and MOCU. This is the \emph{online} execution only; policy training (e.g.\ REINFORCE with terminal MOCU) is done separately.
\item \textbf{\S\ref{sec:gamma_star} (Algorithm 2):} Computes $\gamma^\ast(\vartheta)$ for a given parameter $\vartheta$: the minimal control gain such that the closed-loop system under the reference contingency satisfies frequency and ROCOF limits. Implemented by bracketing (find a safe upper bound) then bisection (narrow to the smallest safe $\gamma$).
\item \textbf{\S\ref{sec:limitations}:} States limitations (myopic vs terminal, parametrization, observation model, bisection validity).
\end{itemize}

% ============================================================
% Section 1: Framework — goals, model, and sequential design
% ============================================================
\section{Framework: Goals, Model, and Sequential Bayesian Design}
\label{sec:setup}

This section defines every object and assumption needed to run the two algorithms: the physical goal and stability condition, the swing-equation model, latent parameters and prior, design space, observation (ROCOF) and likelihood, posterior update (with numerical stability), decision cost and MOCU, and the sequential optimization goal. We use the same sequential design structure as in Foster et al.\ (2021): latent $\vartheta$, design $\xi$, observation $y$, belief update, and inner/outer sampling; our objective is expected terminal MOCU.

\noindent\textit{Notation.} Angles in rad, time in s, $f$ in Hz. Subscripts $i,j$ index buses; $t=0,1,\dots,T$ indexes steps (and beliefs: $p_0$ is the prior, $p_t$ is the belief after step $t$). \textbf{Discrete belief (inner):} the belief is a distribution over a finite set of parameter values (grid or particles). The \emph{inner index} $\ninner=1,\dots,\Ninner$ runs over these $\Ninner$ points (in code: grid size or \texttt{num\_inner\_samples}/$\Ninner$). The $\ninner$-th parameter value is $\vartheta_{\ninner}$ and its probability weight at step $t$ is $p_t^{\ninner}$; thus $\sum_{\ninner=1}^{\Ninner} p_t^{\ninner}=1$. So $\Ninner$ is the \textbf{inner} dimension. \textbf{Inner/outer (DAD convention):} one \emph{episode} is one run over steps $t=0,\dots,T$. Within each step, to evaluate the greedy objective $R(\xi)=\mathbb{E}_y[\mathrm{MOCU}(p_t(\xi,y))]$: \textbf{outer loop} ($\Nouter$): simulate $\Nouter$ possible future observations (sample $\vartheta^{(\nouter)}\sim p_{t-1}$, $y^{(\nouter)}\sim p(y\mid\xi,\vartheta^{(\nouter)})$ for $\nouter=1,\dots,\Nouter$; in code \texttt{num\_outer\_samples}). \textbf{Inner loop} ($\Ninner$): for each such observation, update the belief over the parameter grid---compute likelihood and posterior over the $\Ninner$ points using the \emph{inner index} $\ninner=1,\dots,\Ninner$ (in code \texttt{num\_inner\_samples} or grid size $\Ninner$). So: outer = simulating futures; inner = belief update over the $\Ninner$ points (Foster et al., 2021; \texttt{dad-main}).

\medskip
\noindent
\textbf{Key symbols (quick reference).} $\vartheta$ = latent parameter (e.g.\ $(M,K)$); $\xi=(b,A,T_p)$ = design (probe bus, amplitude, duration); $y$ = scalar observation (max ROCOF + noise); $p_t$ = belief (posterior over $\vartheta$) after step $t$; $h_t$ = history of designs and observations up to step $t$; $\Map(\vartheta,\xi)$ = deterministic forward map (max ROCOF from ODE); $\gamma^\ast(\vartheta)$ = minimal safe gain for parameter $\vartheta$; $\hat{\gamma}(p)$ = Bayes-optimal deployed gain under belief $p$ (weighted median of $\gamma^\ast$); $J(\gamma,\vartheta)=|\gamma-\gamma^\ast(\vartheta)|$ = operational cost; $\mathrm{MOCU}(p)$ = mean objective cost of uncertainty = $\mathbb{E}_{\vartheta\sim p}[J(\hat{\gamma}(p),\vartheta)]$; $\pi_\phi(\cdot\mid h)$ = amortized (neural) policy that outputs the next design given history; $\xi_{\mathrm{ref}}$ = reference probe used only when evaluating $\gamma^\ast$.

% -----------------------------------------------------------------------
\subsection{High-level goal and stability condition}
\label{sec:definitions}

\paragraph{High-level goal.}
The \emph{high-level goal} is to keep the power system \emph{stable} (secure) under a reference contingency. We learn the unknown system parameters via sequential experiments so that we can deploy a control gain that satisfies grid security limits. Stability is defined by \emph{operational security constraints} (grid codes, e.g.\ ENTSO-E, NERC; Kundur, 1994).

\paragraph{Stability condition.}
Over the contingency horizon $[0,T_c]$, the system frequency $f(t)$ and its rate of change (ROCOF) $\dot f(t)$ must stay within limits. The \emph{stability condition} is
\[
\boxed{\;\forall t\in[0,T_c],\quad |\dot f(t)|\le r_{\max}\ \land\ f(t)\ge f_{\min}\;},
\]
where $f=f_{\mathrm{nom}}+\omega/(2\pi)$ and $|\dot f|=|\dot\omega|/(2\pi)$ (e.g.\ at a reference bus or worst over buses). So ``stable'' means: (1) ROCOF stays below $r_{\max}$ (e.g.\ 0.1\,Hz/s; grid codes limit rate of change of frequency), and (2) frequency stays above $f_{\min}$ (e.g.\ 49.8\,Hz at 50\,Hz nominal; avoids underfrequency load shedding). We specify only a \emph{lower} bound $f_{\min}$ on frequency (and not an upper bound $f_{\max}$) because the reference contingency we consider (e.g.\ a probe or loss-of-generation event) primarily causes frequency to \emph{drop}; the binding constraint is then the frequency \emph{nadir} $f\ge f_{\min}$ (to avoid underfrequency load shedding and instability). Grid codes often also impose $f\le f_{\max}$ (overfrequency); the formulation can be extended to include it if the contingency can drive frequency upward. The \emph{minimal safe gain} $\gamma^\ast(\vartheta)$ is the smallest control gain $\gamma\ge 0$ such that, under the reference contingency and control $u^{\mathrm{ctrl}}=\gamma g_i\omega_i$, this condition holds. Deploying $\gamma\ge\gamma^\ast(\vartheta_{\mathrm{true}})$ keeps the system stable for the true parameter; the purpose of learning is to infer $\vartheta$ and choose such a gain.

% -----------------------------------------------------------------------
\subsection{Dynamics model}
\label{sec:model}

\paragraph{Network.}
The power network is $\mathcal{G}=(\mathcal{V},\mathcal{E},\bm{B})$: $|\mathcal{V}|=N$ buses (e.g.\ IEEE-14, $N=14$), $\mathcal{E}$ edges, $\bm{B}=(B_{ij})$ susceptance matrix. Other fixed data: mechanical power $P_{m,i}$, damping $D_i$, control shape $\{g_i\}$, sampling rate $f_s$, observation window $T_{\mathrm{obs}}$, contingency horizon $T_c$, security limits $(r_{\max},f_{\min})$, and reference probe $\xi_{\mathrm{ref}}=(b_{\mathrm{ref}},A_{\mathrm{ref}},T_{\mathrm{ref}})$.

\paragraph{Swing equation (second-order Kuramoto).}
For each bus $i\in\mathcal{V}$:
\[
\dot{\theta}_i=\omega_i,
\qquad
M_i\,\dot{\omega}_i=P_{m,i}-P_{e,i}(\bm{\theta})-(D_i+K_i)\omega_i+u_i,
\]
with $P_{e,i}(\bm{\theta})=\sum_{j\in\mathcal{N}_i}B_{ij}\sin(\theta_i-\theta_j)$. Injections are $u_i=u_i^{\mathrm{probe}}+u_i^{\mathrm{ctrl}}$. \textbf{Probe} $u^{\mathrm{probe}}$: Hann-window at a chosen bus; in the sBOED loop it is set by design $\xi_t=(b_t,A_t,T_p)$; for $\gamma^\ast$ evaluation we use the fixed \emph{reference probe} $\xi_{\mathrm{ref}}$ (reference contingency). \textbf{Control} $u^{\mathrm{ctrl}}_i=\gamma\,g_i\,\omega_i$ (proportional frequency feedback; $\gamma$ is the gain). During the observation step we set $\gamma=0$ so the response is due only to the probe. The Hann window at bus $b$ is $u^{\mathrm{probe}}_b(t;\xi)=A\,\frac{1-\cos(2\pi t/T_p)}{2}\,\mathbf{1}_{[0,T_p]}(t)$, $u^{\mathrm{probe}}_{i\neq b}=0$.

% -----------------------------------------------------------------------
\subsection{Latent parameters and latent space}
\label{sec:latent}

\paragraph{Latent parameters.}
The unknown parameters are \emph{per-bus} inertia $M_i$ and damping gain $K_i$ ($i=1,\dots,N$).

\paragraph{Latent space.}
The \emph{latent space} is $\Theta\subseteq\mathbb{R}_+^{2N}$ with
\[
\vartheta=(M_1,\dots,M_N,K_1,\dots,K_N)\in\Theta.
\]
The true parameter $\vartheta_{\mathrm{true}}\in\Theta$ is fixed but unknown. We represent belief at step $t$ by a discrete set (grid or particles) $\{\vartheta_{\ninner}\}_{\ninner=1}^{\Ninner}$ with weights $\{p_t^{\ninner}\}_{\ninner=1}^{\Ninner}$, $\sum_{\ninner=1}^{\Ninner} p_t^{\ninner}=1$. For large $N$, a full $2N$-dimensional discretization is infeasible; implementations may use a reduced parametrization or particle filters (see \S\ref{sec:limitations}).

% -----------------------------------------------------------------------
\subsection{Prior}
\label{sec:prior}

\paragraph{Prior over $\Theta$.}
The prior $p_0(\vartheta)$ is set over $\Theta$. Typically we take a uniform prior over a bounded box (e.g.\ $M_i\in[M_{\min},M_{\max}]$, $K_i\in[K_{\min},K_{\max}]$) or over a discrete set $\{\vartheta_{\ninner}\}_{\ninner=1}^{\Ninner}$ with initial weights $\{p_0^{\ninner}\}_{\ninner=1}^{\Ninner}$ satisfying $\sum_{\ninner=1}^{\Ninner} p_0^{\ninner}=1$. The discrete set can be a grid over a reduced parametrization or a particle sample; the formulas below hold for any such representation.

% -----------------------------------------------------------------------
\subsection{Design and design space (order matters)}
\label{sec:design}

\paragraph{Design space.}
The \emph{design space} $\Xi$ is the set of single-step probe designs. Each design is $\xi=(b,A,T_p)\in\Xi$: bus $b\in\{1,\dots,N\}$, amplitude $A$, duration $T_p$. At step $t$ we choose $\xi_t\in\Xi$. \textbf{Do not confuse:} the \emph{reference probe} $\xi_{\mathrm{ref}}$ is a \emph{fixed} scenario used only when computing $\gamma^\ast(\vartheta)$ (Algorithm~\ref{alg:gamma_star}); the \emph{design} $\xi_t$ is the probe we \emph{choose} at step $t$ in the sequential loop (Algorithm~\ref{alg:dad_mocu_loop}) to collect observation $y_t$.

\paragraph{Order matters.}
A \emph{design sequence} is an ordered list $\xi_{1:T}=(\xi_1,\dots,\xi_T)$. We assume each probe is run from fixed initial conditions, so the deterministic $\Map(\vartheta,\xi)$ does not depend on the order of previous probes; order still matters for the decision problem because different sequences yield different observation noise and adaptive choices, hence different posteriors and terminal MOCU. With $|\Xi|$ designs and without replacement, there are $|\Xi|(|\Xi|-1)\cdots(|\Xi|-T+1)$ sequences (e.g.\ $42\times 41\times 40\times 39$ for $|\Xi|=42$, $T=4$)---exhaustive search is infeasible. The myopic rule minimizes only one-step-ahead expected MOCU; because the marginal value of a design can depend on the set already chosen (complementarity), myopic is generally suboptimal for terminal MOCU and fails to plan for future information gain.

% -----------------------------------------------------------------------
\subsection{Observation and why ROCOF only}
\label{sec:observation}

\paragraph{Forward map.}
Let $f_s$ be the sampling rate and $\Delta t=1/f_s$. From the ODE solution (initial conditions fixed, e.g.\ equilibrium), define frequency deviation $\Delta f_i(t)=\omega_i(t)/(2\pi)$ at bus $i$, and its discrete-time version $\Delta f_i[k]=\Delta f_i(k\Delta t)$ at sample index $k$. The \emph{rate of change of frequency} (ROCOF) at bus $i$ and time sample $k$ is $\mathrm{ROCOF}_i[k]=(\Delta f_i[k]-\Delta f_i[k-1])/\Delta t$ (units Hz/s). Here $i\in\{1,\dots,N\}$ is the \textbf{bus index} and $k$ is the \textbf{discrete time-sample index} over the observation window. The deterministic \emph{forward map} is the maximum absolute ROCOF over all buses and time samples:
\[
\Map(\vartheta,\xi)=\max_{i,k}\left|\mathrm{ROCOF}_i[k]\right|.
\]

\paragraph{Why ROCOF as observation (and references).}
ROCOF is constrained by grid codes (ENTSO-E; NERC BAL-001-TRE-2; Kundur, 1994) and is physically linked to system inertia: after a power imbalance, the rate of change of frequency is inversely related to effective inertia (lower inertia typically yields higher ROCOF). Using ROCOF (or its maximum over buses and time) as the observable is therefore standard in frequency security and inertia assessment (e.g.\ Kundur, 1994; Ulbig et al., 2014; Australian Energy Market Operator, 2020). We observe a single scalar $y$ per experiment. Using the \emph{maximum ROCOF} over buses and time is (i) \emph{decision-relevant}: it directly relates to the stability condition $|\dot f|\le r_{\max}$; (ii) \emph{informative} about inertia and damping; (iii) \emph{tractable}: a scalar observation gives a simple Gaussian likelihood and keeps the design loop computationally feasible (Foster et al., 2021). Richer observations (e.g.\ full $\omega_i(t)$ or multiple features) could be used at higher cost.

\paragraph{Observation model and $\sigma_{\mathrm{feat}}$.}
The observation is $y=\Map(\vartheta,\xi)+\epsilon$, $\epsilon\sim\mathcal{N}(0,\sigma_{\mathrm{feat}}^2)$. The term $\sigma_{\mathrm{feat}}$ is the \emph{observation noise standard deviation} (in Hz/s) at the \textbf{scalar-feature} level: the measured $y$ is the true max ROCOF plus independent Gaussian noise with mean zero and variance $\sigma_{\mathrm{feat}}^2$. Noise is modeled at the feature level; sensor-level noise before taking the max would yield a non-Gaussian distribution for the max, but we use the Gaussian model for tractability (standard in OED when the forward map is expensive).

% -----------------------------------------------------------------------
\subsection{Likelihood}
\label{sec:likelihood}

\paragraph{Likelihood calculation.}
The \emph{likelihood} $p(y\mid\vartheta,\xi)$ is the density of $y$ given parameter $\vartheta$ and design $\xi$. With the model above it is Gaussian:
\[
p(y\mid \vartheta,\xi)=\frac{1}{\sqrt{2\pi\sigma_{\mathrm{feat}}^2}}
\exp\left(-\frac{\big(y-\Map(\vartheta,\xi)\big)^2}{2\sigma_{\mathrm{feat}}^2}\right).
\]
On the discrete set $\{\vartheta_{\ninner}\}_{\ninner=1}^{\Ninner}$ (inner index $\ninner$), we write $\mathcal{L}_{\ninner}^{(t)}=p(y_t\mid\vartheta_{\ninner},\xi_t)$ with mean $\mu_{\ninner}^{(t)}=\Map(\vartheta_{\ninner},\xi_t)$ and variance $\sigma_{\mathrm{feat}}^2$.

% -----------------------------------------------------------------------
\subsection{Posterior update and numerical stability}
\label{sec:posterior}

\paragraph{Exact discrete Bayesian update.}
At step $t$, the belief \emph{before} the observation is $p_{t-1}(\vartheta)$. Upon applying design $\xi_t$ and observing the scalar feature $y_t$, the belief is updated via Bayes' rule: $p_t(\vartheta)\propto p_{t-1}(\vartheta)\,p(y_t\mid\vartheta,\xi_t)$. Over the discrete support $\{\vartheta_{\ninner}\}_{\ninner=1}^{\Ninner}$, the theoretical update scales the prior probability mass of each candidate point by its corresponding likelihood $\mathcal{L}_{\ninner}^{(t)} = p(y_t\mid\vartheta_{\ninner},\xi_t)$. The unnormalized posterior weights and the normalization constant are given by:
\[
\tilde p_t^{\ninner} = p_{t-1}^{\ninner} \cdot \mathcal{L}_{\ninner}^{(t)}, \qquad Z_t = \sum_{\ninner=1}^{\Ninner} \tilde p_t^{\ninner}, \qquad p_t^{\ninner} = \frac{\tilde p_t^{\ninner}}{Z_t}.
\]
Here $Z_t$ is the \textbf{normalization constant} (also called the \emph{evidence} or marginal likelihood of $y_t$ given $\xi_t$ and $p_{t-1}$); it ensures $\sum_{\ninner=1}^{\Ninner} p_t^{\ninner}=1$.

\paragraph{Log-domain implementation (Log-Sum-Exp).}
In practice, a direct implementation of the multiplicative update is susceptible to severe numerical underflow. Because the Gaussian likelihood decays exponentially with the squared error $\big(y_t-\Map(\vartheta_{\ninner},\xi_t)\big)^2$, evaluating $\mathcal{L}_{\ninner}^{(t)}$ for highly improbable parameter candidates yields values that exceed the limits of standard floating-point precision, rounding to absolute zero. If all $\tilde p_t^{\ninner}$ underflow, $Z_t$ evaluates to zero, leading to an undefined posterior. To guarantee numerical stability, the update must be executed entirely in the log-domain. The unnormalized log-posterior is formulated as:
\[
\log \tilde p_t^{\ninner} = \log p_{t-1}^{\ninner} - \frac{1}{2}\log(2\pi\sigma_{\mathrm{feat}}^2) - \frac{\big(y_t - \Map(\vartheta_{\ninner}, \xi_t)\big)^2}{2\sigma_{\mathrm{feat}}^2}.
\]
To compute the normalized weights $p_t^{\ninner}=\exp(\log \tilde p_t^{\ninner} - \log Z_t)$ without overflow or underflow, we employ the log-sum-exp trick. By defining a shift constant $c_t = \max_{\ninner} \log \tilde p_t^{\ninner}$, the log-evidence is computed as:
\[
\log Z_t = c_t + \log \sum_{\ninner'=1}^{\Ninner} \exp\big(\log \tilde p_t^{\ninner'} - c_t\big).
\]
The final normalized posterior probabilities are then accurately recovered via $p_t^{\ninner} = \exp\big(\log \tilde p_t^{\ninner} - \log Z_t\big)$ (Chopin \& Papaspiliopoulos, 2020).

% -----------------------------------------------------------------------
\subsection{Decision objective, cost function, and MOCU}
\label{sec:cost_mocu}

\paragraph{Minimal safe gain and cost function.}
Under the reference probe $\xi_{\mathrm{ref}}$ and constraints $(r_{\max},f_{\min})$ over $[0,T_c]$, the \emph{minimal safe gain} is
\[
\gamma^{\ast}(\vartheta)=\inf\left\{\gamma\ge 0 \;:\; \text{trajectory with } u^{\mathrm{probe}}(\cdot;\xi_{\mathrm{ref}}),\; u^{\mathrm{ctrl}}=\gamma g_i\omega_i \text{ satisfies } |\dot f|\le r_{\max},\; f\ge f_{\min}\right\}
\]
(computed by Algorithm~\ref{alg:gamma_star}). The \emph{operational cost} (loss when deploying $\gamma$ and truth is $\vartheta$) is
\[
J(\gamma,\vartheta)=\left|\gamma-\gamma^\ast(\vartheta)\right|
\]
(Boluki et al., 2018; Imani et al., 2018). Given belief $p$, the \emph{Bayes-optimal deployed gain} is the weighted median (Ferguson, 1967):
\[
\hat{\gamma}(p)=\argmin_{\gamma\ge 0}\mathbb{E}_{\vartheta\sim p}[J(\gamma,\vartheta)]
=\wmed\!\left(\{\gamma^\ast(\vartheta_{\ninner})\}_{\ninner=1}^{\Ninner},\{p^{\ninner}\}_{\ninner=1}^{\Ninner}\right),
\]
where $\wmed$ is the value $\gamma$ such that the $p$-weighted mass on $\{\vartheta_{\ninner}:\gamma^\ast(\vartheta_{\ninner})\le\gamma\}$ is at least $1/2$.

\paragraph{MOCU definition and formula.}
The \emph{mean objective cost of uncertainty} (MOCU) is the expected cost when deploying the Bayes-optimal gain under belief $p$ (Boluki et al., 2018; Imani et al., 2018):
\[
\MOCU(p)=\mathbb{E}_{\vartheta\sim p}\left[J(\hat{\gamma}(p),\vartheta)\right]
=\sum_{\ninner=1}^{\Ninner} p^{\ninner}\left|\gamma^\ast_{\ninner}-\hat{\gamma}(p)\right|,
\]
with $\gamma^\ast_{\ninner}=\gamma^\ast(\vartheta_{\ninner})$. \textbf{In plain language:} MOCU is the expected absolute gap between the optimal gain for the true $\vartheta$ and the gain we deploy; lower MOCU means better decision quality (less uncertainty that matters for the deployment decision). In a single run the index $t$ runs from $0$ to $T$: $t=0$ is the \emph{initial} belief $p_0$ (prior; no designs applied), and \emph{all methods have the same MOCU at $t=0$}; designs are applied at steps $t=1,2,\dots,T$, and after step $t$ the belief is $p_t$. \emph{Initial MOCU in experiments} is $\mathrm{MOCU}(p_0)$: MOCU under the prior. With a discrete belief it is $\sum_{\ninner=1}^{\Ninner} p_0^{\ninner}\,|\gamma^\ast_{\ninner}-\hat{\gamma}(p_0)|$ with $\hat{\gamma}(p_0)=\wmed(\{\gamma^\ast_{\ninner}\}_{\ninner=1}^{\Ninner},\{p_0^{\ninner}\}_{\ninner=1}^{\Ninner})$. In our implementation the prior is uniform over the initial uncertainty box $[M_{\mathrm{lower}},M_{\mathrm{upper}}]\times[K_{\mathrm{lower}},K_{\mathrm{upper}}]$; we sample $(M,K)$ from that box, compute $\gamma^\ast(M,K)$ for each sample under the reference contingency, set $\hat{\gamma}=\mathrm{median}(\gamma^\ast)$ and $\mathrm{MOCU}=\mathrm{mean}(|\gamma^\ast-\hat{\gamma}|)$ (same formula, Monte Carlo over the continuous prior). Each new observation adds zero or positive information, so $\mathrm{MOCU}(p_t)$ is \emph{non-increasing} in $t$ along any given run. \emph{Order matters}: different design sequences yield different posteriors and thus different (non-increasing) MOCU curves and different \emph{terminal} MOCU $\mathrm{MOCU}(p_T)$; the myopic rule is \emph{generally suboptimal} for terminal MOCU, since it optimizes only the next step and does not account for complementarity. The goal is to choose a sequence that achieves low terminal MOCU.

% -----------------------------------------------------------------------
\subsection{Optimization goal (sequential design)}
\label{sec:opt_goal}

\paragraph{Optimization goal.}
The \emph{optimization goal} of the sequential design is to minimize \emph{expected terminal MOCU}:
\[
\min_{\xi_1,\dots,\xi_T}\;\mathbb{E}\big[\mathrm{MOCU}(p_T)\big],
\]
where the expectation is over the randomness in observations $y_1,\dots,y_T$ (and optionally $\vartheta_{\mathrm{true}}\sim p_0$), and $p_T$ is the belief after $T$ steps. \textbf{In plain language:} we choose a design \emph{policy} (or sequence) so that at the end of $T$ experiments the expected remaining cost of uncertainty is as small as possible; then we deploy $\hat{\gamma}(p_T)$. Per-step MOCU is computed for monitoring and for greedy heuristics but is \emph{not} the quantity being optimized.

% -----------------------------------------------------------------------
\subsection{Assumptions and inputs/outputs}
\label{sec:assumptions_io}

\paragraph{Assumptions.}
We assume the safety predicate $\Safe(\gamma;\vartheta)$ is non-decreasing in $\gamma$ (safe set $[\gamma^\ast(\vartheta),\infty)$ so bisection is valid) and that observation noise is Gaussian. The discrete set $\{\vartheta_{\ninner}\}_{\ninner=1}^{\Ninner}$ (inner index $\ninner$) is such that $\gamma^\ast(\vartheta_{\ninner})$ can be evaluated for each $\ninner=1,\dots,\Ninner$. Initial conditions for the swing ODE are fixed so that $\Map(\vartheta,\xi)$ is well-defined.

\paragraph{Inputs and outputs.}
Inputs to Algorithm~\ref{alg:dad_mocu_loop}: $\mathcal{G}$; discrete set $\{\vartheta_{\ninner}\}_{\ninner=1}^{\Ninner}$ and initial belief $\{p_0^{\ninner}\}_{\ninner=1}^{\Ninner}$; design space $\Xi$; $\sigma_{\mathrm{feat}}$, $f_s$, $T_{\mathrm{obs}}$; $(r_{\max},f_{\min})$, $T_c$, $\xi_{\mathrm{ref}}$, $\{g_i\}$, horizon $T$. Outputs: beliefs $\{p_t\}_{t=0}^{T}$, gains $\{\hat{\gamma}(p_t)\}_{t=0}^{T}$, MOCU at $t=0,\dots,T$ (with $t=0$ initial, $t=T$ terminal). At $t=0$ all methods have the same MOCU; the \emph{ultimate} quantity of interest is terminal MOCU at $t=T$, i.e.\ $\mathrm{MOCU}(p_T)$.

\newpage
% ============================================================
% Section 2: Sequential design and learning
% ============================================================
\section{Sequential design and learning}
\label{sec:sequential}

This section gives \textbf{Algorithm~\ref{alg:dad_mocu_loop}}: the \emph{online} execution of one episode of sequential design. The goal is to minimize expected \emph{terminal} MOCU $\mathbb{E}[\mathrm{MOCU}(p_T)]$. The algorithm assumes a \textbf{pre-trained} amortized policy $\pi_\phi$ (e.g.\ trained by REINFORCE with terminal MOCU as reward; training is separate and not shown here). In one episode: at $t=0$ belief is $p_0$ (prior); at each step $t=1,\dots,T$ we \emph{sample} design $\xi_t$ from $\pi_\phi(\cdot\mid h_{t-1})$, observe $y_t$, update belief to $p_t$ (log-domain), and compute $\hat{\gamma}(p_t)$ and $\mathrm{MOCU}(p_t)$. Alternative: a \emph{greedy} (myopic) rule would minimize one-step-ahead expected MOCU at each step (outer loop over simulated futures, inner loop over belief points); that is tractable but generally suboptimal for terminal MOCU. \textbf{Inner/outer (for greedy only):} outer = $\Nouter$ simulated future observations; inner = belief update over $\Ninner$ points. With the amortized policy, we do \emph{not} run outer/inner at deployment---we only run the policy and the belief update for the single real observation per step.

\begin{algorithm}[H]
\caption{Sequential Bayesian Experimental Design (sBOED) using Amortized Policy}
\label{alg:dad_mocu_loop}
\small
\begin{algorithmic}[1]
\Require Network $\mathcal{G}$; discrete support $\{\vartheta_{\ninner}\}_{\ninner=1}^{\Ninner}$ and prior $\{p_0^{\ninner}\}_{\ninner=1}^{\Ninner}$; noise $\sigma_{\mathrm{feat}}$; precomputed $\gamma^\ast_{\ninner}=\gamma^\ast(\vartheta_{\ninner})$ for all $\ninner$; horizon $T$; \textbf{pre-trained} amortized policy $\pi_\phi$ (e.g.\ neural network trained offline to minimize expected terminal MOCU).
\Ensure One episode: beliefs $\{p_t\}_{t=0}^{T}$, optimal gains $\{\hat{\gamma}(p_t)\}_{t=0}^{T}$, and MOCU trajectory. The \emph{terminal} output of interest is $\mathrm{MOCU}(p_T)$.

\State $h_0 \gets \emptyset$
\State $\hat{\gamma}(p_0) \gets \wmed\big(\{\gamma^\ast_{\ninner}\}_{\ninner=1}^{\Ninner}, \{p_0^{\ninner}\}_{\ninner=1}^{\Ninner}\big)$
\State $\MOCU(p_0) \gets \sum_{\ninner=1}^{\Ninner} p_0^{\ninner}\,|\gamma^\ast_{\ninner}-\hat{\gamma}(p_0)|$

\For{$t=1,\dots,T$}
    \State \Comment{\textbf{1. Design Phase}}
    \State $\xi_t \sim \pi_\phi(\cdot \mid h_{t-1})$ \Comment{Sample probe from trained neural network policy}

    \State \Comment{\textbf{2. Physical Observation}}
    \State $y_t \gets \Map(\vartheta_{\mathrm{true}}, \xi_t) + \epsilon_t$, where $\epsilon_t \sim \mathcal{N}(0, \sigma_{\mathrm{feat}}^2)$

    \State \Comment{\textbf{3. Exact Discrete Belief Update (Log-domain)}}
    \For{$\ninner=1,\dots,\Ninner$}
        \State $\log \tilde{p}_t^{\ninner} \gets \log p_{t-1}^{\ninner} - \frac{1}{2}\log(2\pi\sigma_{\mathrm{feat}}^2) - \frac{\big(y_t - \Map(\vartheta_{\ninner}, \xi_t)\big)^2}{2\sigma_{\mathrm{feat}}^2}$
    \EndFor
    \State $c_t \gets \max_{\ninner} \log \tilde{p}_t^{\ninner}$ \Comment{Shift constant for Log-Sum-Exp}
    \State $\log Z_t \gets c_t + \log \sum_{\ninner'=1}^{\Ninner} \exp\big(\log \tilde{p}_t^{\ninner'} - c_t\big)$
    \For{$\ninner=1,\dots,\Ninner$}
        \State $p_t^{\ninner} \gets \exp\big(\log \tilde{p}_t^{\ninner} - \log Z_t\big)$
    \EndFor

    \State \Comment{\textbf{4. Deploy Optimal Gain \& Track MOCU}}
    \State $\hat{\gamma}(p_t) \gets \wmed\big(\{\gamma^\ast_{\ninner}\}_{\ninner=1}^{\Ninner}, \{p_t^{\ninner}\}_{\ninner=1}^{\Ninner}\big)$
    \State $\MOCU(p_t) \gets \sum_{\ninner=1}^{\Ninner} p_t^{\ninner}\,|\gamma^\ast_{\ninner}-\hat{\gamma}(p_t)|$
    \State $h_t \gets h_{t-1} \cup \{(\xi_t, y_t)\}$
\EndFor
\State \Return $\{p_t\}_{t=0}^{T}$, $\{\hat{\gamma}(p_t)\}_{t=0}^{T}$, $\{\MOCU(p_t)\}_{t=0}^{T}$
\end{algorithmic}
\end{algorithm}

% ============================================================
% Section 3: Minimal safe gain
% ============================================================
\section{Computing the minimal safe gain}
\label{sec:gamma_star}

This section gives \textbf{Algorithm~\ref{alg:gamma_star}}: for a \emph{given} parameter $\vartheta$, compute the \emph{minimal safe gain} $\gamma^\ast(\vartheta)$---the smallest $\gamma\ge 0$ such that the closed-loop trajectory under the \textbf{reference probe} $\xi_{\mathrm{ref}}$ and control $u^{\mathrm{ctrl}}=\gamma g_i\omega_i$ satisfies $|\dot f|\le r_{\max}$ and $f\ge f_{\min}$ on $[0,T_c]$. This is used to evaluate $\gamma^\ast_{\ninner}$ for each belief point $\vartheta_{\ninner}$ (precomputed or on the fly). We assume the set of safe gains is an interval $[\gamma^\ast(\vartheta),\infty)$ (monotonicity in $\gamma$; see \S\ref{sec:limitations}), so \textbf{bisection} is valid: first \emph{bracket} to find an upper bound $\gamma_U$ that is safe, then \emph{bisect} to narrow to the smallest such $\gamma$. If no safe $\gamma$ exists below a hard cap $\gamma_{\mathrm{cap}}$, the algorithm returns $\gamma_{\mathrm{cap}}$ (``unsolvable'' in the sense that the system cannot be stabilized within the cap).

\begin{algorithm}[H]
\caption{Minimal Safe Gain $\gamma^\ast(\vartheta)$}
\label{alg:gamma_star}
\small
\begin{algorithmic}[1]
\Require Parameter $\vartheta$; network $\mathcal{G}$; limits $(r_{\max}, f_{\min})$; horizon $T_c$; reference probe $\xi_{\mathrm{ref}}$; tolerances $\varepsilon$ (bisection stop), $N_{\mathrm{expand}}$ (bracketing iterations); caps $\gamma_{\max}$ (initial upper guess), $\gamma_{\mathrm{cap}}$ (hard cap; return this if no safe $\gamma$ below it).
\Ensure Minimal safe gain $\gamma^\ast(\vartheta)$; or $\gamma_{\mathrm{cap}}$ if the system cannot be stabilized (no $\gamma\le\gamma_{\mathrm{cap}}$ satisfies the safety condition).

\State \textbf{Define} $\Safe(\gamma;\vartheta) = \text{True}$ \textbf{iff} ODE trajectory under $\xi_{\mathrm{ref}}$ and $u^{\mathrm{ctrl}}=\gamma g_i\omega_i$ satisfies $|\dot{f}| \le r_{\max}$ and $f \ge f_{\min}$ on $[0,T_c]$.

\State \Comment{\textbf{1. Bracketing Phase}}
\State $\gamma_L \gets 0$
\State $\gamma_U \gets \gamma_{\max}$
\For{$k=1,\dots,N_{\mathrm{expand}}$}
    \If{$\Safe(\gamma_U;\vartheta)$ is True}
        \State \textbf{break} \Comment{Valid upper bound found}
    \EndIf
    \State $\gamma_U \gets \min(2\gamma_U, \gamma_{\mathrm{cap}})$
    \If{$\gamma_U = \gamma_{\mathrm{cap}}$}
        \State \textbf{break} \Comment{Hit hard cap}
    \EndIf
\EndFor

\If{\textbf{not} $\Safe(\gamma_U;\vartheta)$}
    \State \Return $\gamma_{\mathrm{cap}}$ \Comment{System cannot be stabilized}
\EndIf

\State \Comment{\textbf{2. Bisection Phase}}
\While{$\gamma_U - \gamma_L > \varepsilon$}
    \State $\gamma_M \gets (\gamma_L + \gamma_U) / 2$
    \If{$\Safe(\gamma_M;\vartheta)$ is True}
        \State $\gamma_U \gets \gamma_M$ \Comment{Midpoint is safe; lower the upper bound}
    \Else
        \State $\gamma_L \gets \gamma_M$ \Comment{Midpoint is unsafe; raise the lower bound}
    \EndIf
\EndWhile

\State \Return $\gamma_U$
\end{algorithmic}
\end{algorithm}

% ============================================================
% Section 4: Limitations and modeling choices
% ============================================================
\section{Limitations and modeling choices}
\label{sec:limitations}

This section states limitations and modeling choices explicitly so that the scope of the framework is clear; each subsection indicates what is assumed and possible extensions.

\subsection{Sequential design and myopic heuristic}

The objective is expected \emph{terminal} MOCU $\mathbb{E}[\mathrm{MOCU}(p_T)]$; the sequential design is order-dependent and MOCU is non-increasing in $t$. The myopic rule minimizes only one-step-ahead expected MOCU and is \emph{generally suboptimal} for terminal MOCU (rarely optimal; in edge cases where designs provide independent information, greedy can align with the optimum); it is a tractable heuristic. In the current setup (scalar ROCOF, reset each probe, reduced $\vartheta$, finite $\Xi$), complementarity between designs may be weak, so myopic may perform close to optimal in practice; to expose a larger gap one could use richer observations, higher-dimensional $\vartheta$, larger $T$, or terminal-aware baselines where tractable.

\subsection{Parametrization and tractability}

The latent parameter $\vartheta=(M_1,\dots,M_N,K_1,\dots,K_N)\in\mathbb{R}_+^{2N}$ makes a full discrete belief over $2N$ dimensions infeasible for large $N$. Implementations may use a reduced parametrization (e.g.\ common $M$, $K$ and a small grid) or Sequential Monte Carlo (Chopin \& Papaspiliopoulos, 2020; Drovandi et al., 2013). The algorithms in this document apply to any belief representation (grid or particles).

\subsection{Observation model and noise}

Observations are scalar ($y=\Map(\vartheta,\xi)+\epsilon$, $\epsilon\sim\mathcal{N}(0,\sigma_{\mathrm{feat}}^2)$). The maximum of noisy signals is not strictly Gaussian; we use a Gaussian likelihood for tractability. During observation we set $u^{\mathrm{ctrl}}\equiv 0$, which isolates parameter identification but is a security concern in operation; a closed-loop formulation with baseline gain $\gamma_{\mathrm{base}}>0$ during probing would extend the model for deployment.

\subsection{Bisection validity and deployment}

Algorithm~\ref{alg:gamma_star} assumes the safe set is $[\gamma^\ast,\infty)$ (safety predicate non-decreasing in $\gamma$). Very high gain can cause oscillations or instability in real systems, so monotonicity may fail; we assume it holds over the operational range and reference contingency used here. If not, $\gamma^\ast$ would require grid search or constrained optimization instead of bisection.

% ============================================================
% References (APA 7th edition)
% ============================================================
\section*{References}
\addcontentsline{toc}{section}{References}

\begingroup
\raggedright
\setlength{\parindent}{0pt}
\setlength{\leftskip}{0pt}
\setlength{\parskip}{0.6em}

Boluki, S., Qian, X., \& Dougherty, E.~R. (2018). Experimental design via generalized mean objective cost of uncertainty. \emph{arXiv:1805.01143}.\\
\url{https://arxiv.org/abs/1805.01143}

Chopin, N., \& Papaspiliopoulos, O. (2020). \emph{An introduction to sequential Monte Carlo}. Springer.\\
\url{https://doi.org/10.1007/978-3-030-47845-2}

D\"orfler, F., \& Bullo, F. (2012). Synchronization in complex networks of phase oscillators: A survey. \emph{Automatica}, \emph{48}(6), 1089--1106.\\
\url{https://doi.org/10.1016/j.automatica.2012.02.028}

Drovandi, C.~C., McGree, J.~M., \& Pettitt, A.~N. (2013). Sequential Monte Carlo for Bayesian sequentially designed experiments for discrete data. \emph{Computational Statistics \& Data Analysis}, \emph{57}(1), 320--335.\\
\url{https://doi.org/10.1016/j.csda.2012.06.024}

Ferguson, T.~S. (1967). \emph{Mathematical statistics: A decision theoretic approach}. Academic Press.

Foster, A., Ivanova, D.~R., Malik, I., \& Rainforth, T. (2021). Deep adaptive design: Amortizing sequential Bayesian experimental design. In M.~Meila \& T.~Zhang (Eds.), \emph{Proceedings of the 38th International Conference on Machine Learning} (Vol.~139, pp.~3384--3395). PMLR.\\
\url{https://proceedings.mlr.press/v139/foster21a.html}

Imani, M., Dehghannasiri, R., Braga-Neto, U.~M., \& Dougherty, E.~R. (2018). Sequential experimental design for optimal structural intervention in gene regulatory networks based on the mean objective cost of uncertainty. \emph{Cancer Informatics}, \emph{17}, 1--11.\\
\url{https://doi.org/10.1177/1176935118790247}

Kundur, P. (1994). \emph{Power system stability and control}. McGraw-Hill.

Ulbig, A., Borsche, T.~S., \& Andersson, G. (2014). Impact of low rotational inertia on power system stability and operation. \emph{IFAC Proceedings Volumes}, \emph{47}(3), 7290--7297.\\
\url{https://doi.org/10.3182/20140824-6-ZA-1003.02443}

\medskip
\textit{Standards \& guidelines:} ENTSO-E (European Network of Transmission System Operators for Electricity); NERC; NERC BAL-001-TRE-2 (frequency response); Australian Energy Market Operator (AEMO), \emph{Inertia requirements and ROCOF}, 2020.

\endgroup

\end{document}
